\minitopic{来源并不是现成理由的容器}说“我亲眼看见”并没有结束分析。知觉需要感官接触、注意选择、对象分类和信念形成；每一环都可能受环境、训练和偏见影响。说“我记得”同样可能指保存了原有知识、当前有熟悉感，或重构出一个内容。先验判断也不是从真空出现：主体通常先通过经验获得语言、符号与概念，之后才可能凭理解和推理获得不依赖特定观察的确证。来源名称只标出起点，不能替代完整认识链。

\minitopic{知觉理论解释不同层次}直接实在论强调正常经验让主体接触事实本身；意向论强调经验具有可真可假的表征内容；析取论拒绝把真知觉与不可区分幻觉归入决定认识地位的同一根本种类。它们不仅给出不同形而上学图景，也解释理由的方式不同。比较时要固定问题：是在问经验是什么、主体当时能使用何种理由，还是正常知觉为何比幻觉更容易产生知识。用一个层次的答案直接裁决另一层次，常造成虚假胜负。

\minitopic{注意既能增能也能制造盲点}专家训练使放射科医生看见新手忽略的模式，也可能让医生过早套用熟悉分类；社会偏见则可能使证人根本不注意与预期冲突的细节。应区分经验内容是否改变、注意资源分配是否改变、后续信念解释是否改变。即使无法确定认知渗透是否进入经验本身，主体和制度仍能对注意策略负责：采用双读片、盲评或强制搜索清单，能减少只寻找支持既有假设的倾向。认识能动性因此不仅是相信或不相信，也包括决定看哪里、比较什么和何时重新观察。

\minitopic{记忆知识跨越时间而非复制文件}记忆保存论提醒我们，当前记忆资格依赖最初学习与中间保存链；生成论则指出记忆会整合痕迹、背景知识和后来信息。二者可共同解释为什么遗忘原始理由有时无害，有时危险。常识性内容若由稳定学习和反复使用维持，即使不能复述最初证据，仍可能构成知识；对某人的严重指控若只剩来源不明的熟悉感，来源监控失败便是强击败者。关键不是能否回放原场景，而是内容是否通过可接受的保存与修正过程抵达现在。

\begin{argumentbox}
评价认识来源可依次检查五个节点：输入是否与目标对象适当接触；注意和分类是否保留相关信息；信念是否基于该输入形成；跨时间或跨媒介传递是否保存认识地位；当前是否存在足以击败来源链的信息。五节点比“亲眼所见”“我记得很清楚”更能定位失败。
\end{argumentbox}

\minitopic{内省的权威有范围}主体通常比旁观者更有资格报告“我现在疼”“我正在犹豫”，但这不保证他准确解释心理状态的原因、长期倾向或社会意义。感受本身、对感受的概念分类和关于原因的理论推断需要分开。一个人可以无误地报告不适，却误以为不适来自某种食物；也可能因缺少概念资源而无法恰当描述经验。第一人称权威应被理解为在某些对象和条件下的非对称资格，而不是关于心灵的一般不可错论。

\minitopic{先验与归纳显示理性边界}理解“所有单身者都未婚”似乎足以判断命题，但学习相关词语需要经验，命题是否只是语言约定又取决于意义理论。归纳则从过去和样本走向未来与总体，休谟指出其总体可靠性不能用过去归纳成功无循环地证明\parencite[sec.~4]{hume2000}。两类问题共同说明：理性规范既不能完全由个别经验推出，也不能脱离实践。我们可以评价局部推理、校准和替代模型，却不必假装已经从无争议前提证明全部认识方法。

\minitopic{阶段性结论}经验、记忆、内省与先验不是互不相干的四条管道，而是互相支撑又互相校正的来源网络。经验提供世界约束，记忆保存并重组材料，内省报告主体当前状态，先验推理揭示概念和形式关系。知识产生于这些能力与环境的适当合作，也因注意偏差、来源错配、概念贫乏和击败者而失效。来源论的目标不是宣布哪条管道绝对可靠，而是说明每条链怎样被检查和修复。
